\section{Related Work and Positioning}

The target problem is positioned relative to public work on Artemis
trajectory-correction design, rather than being presented as a replacement for
mission-owned analysis. Woffinden et al. describe a correction-burn placement
formulation that incorporates uncertainty, crew-related timing, and entry
constraints \cite{woffinden2023}. Classical learning research has explored
low-thrust guidance and cislunar transfer design, but those settings differ in
dynamics, propulsion, objectives, and validation structure
\cite{sullivan2021,izzo2021,xie2022}. They are valuable precedents for
surrogate modeling and for retaining infeasible cases, not interchangeable
baselines for a human-rated correction benchmark.

Quantum feature-space methods express a classical input through an encoding
and compare encoded states by a quantum-derived similarity
\cite{schuld2019,havlicek2019}. The method can be used in a kernel learner,
but its predictive value depends on an alignment between the data geometry,
the feature map, the measurement, and the task. Theoretical and numerical work
has made this caveat increasingly explicit: hard-to-simulate circuits do not
automatically give a learning advantage, and the data distribution can make a
classical approximation competitive \cite{huang2021,sweke2025}.

The present work is particularly informed by three benchmarking lessons. First,
kernel bandwidth and encoding scale are part of the effective inductive bias,
not negligible implementation details \cite{shaydulin2022}. Second, a
quantum-kernel study should include appropriately strong classical
counterparts, including random features where relevant
\cite{otten2020,sweke2025}. Third, a comparison should report more than its
best run: model choice, split structure, seed variation, capacity, and
post-outcome choices all affect what a claimed gain means
\cite{bowles2024,schnabel2024,pernot2020}.

This paper does not make an exhaustive literature-review claim. The repository
preserves a bounded extraction set of 23 accepted records used to establish the
initial protocol and a separate archival discovery refresh. The refresh did
not change the experiment, thresholds, figures, or outcomes. Literature is
therefore used here to position the benchmark and justify controls, not to
construct a retrospective theory of success.
